\section{Light --- Lib --- Radix Pass}
%%% generated by the blueprint pipeline from TCSlib.Programs.Light.Lib.RadixPass
\subsection*{Light --- Lib --- Radix Pass: Introduction}

Declarations of the Lean module \texttt{TCSlib.Programs.Light.Lib.RadixPass}.

\subsection*{Light --- Lib --- Radix Pass: Results}

\begin{theorem}[Stable pass matching a memory segment]
\lean{Light.exists_stablePass}
\label{Light.exists_stablePass}
\leanok
\uses{Light.IsStablePass, Light.SegN, Light.exists_sortPos_eq, Light.keyAt, Light.readSeg, Light.sortPos, Light.sortPos_lt}
\difficulty{4}
Let $\mu':\mathbb N\to\mathbb Z$, $dst\in\mathbb N$, and let $\pi$ be a finite list of
natural numbers; let $kd:\mathbb N\to\mathbb N$. Suppose that for each index $j$ in
$\pi$, the memory value at $dst$ plus the sorting position of $j$ by the key
$kd(\pi[j])$ equals $\pi[j]$. Then there is a list $\pi'$ of the same length such that,
at each such sorting position, $\pi'$ contains the corresponding original entry
$\pi[j]$, and the memory segment beginning at $dst$ consists exactly of the entries of
$\pi'$ converted to integers.
\end{theorem}

\begin{definition}[Radix pass number field]
\lean{Light.RadixPass.Num}
\label{Light.RadixPass.Num}
\leanok
The radix-pass number field is an alias for the natural number $0$.
\end{definition}

\begin{definition}[Radix pass permutation field]
\lean{Light.RadixPass.Perm}
\label{Light.RadixPass.Perm}
\leanok
\uses{Light.RadixPass.Num}
The radix-pass permutation field is an alias for the natural number $5$.
\end{definition}

\begin{definition}[Radix pass unused field]
\lean{Light.RadixPass.Unused}
\label{Light.RadixPass.Unused}
\leanok
\uses{Light.RadixPass.Num}
The radix-pass unused field is an alias for the natural number $7$.
\end{definition}

\begin{definition}[Radix pass statement]
\lean{Light.radixPass}
\label{Light.radixPass}
\leanok
\uses{Light.Expr, Light.RadixPass.Num, Light.RadixPass.Perm, Light.RadixPass.Unused, Light.Stmt, Light.v}
The radix pass is the sequential composition of two procedure calls. The first call uses
the number field, permutation field, their sum, and the five supplied key, stride,
offset, bound, and counter expressions as arguments; the second uses the sum of the
permutation and number fields, the permutation field, and the number field.
\end{definition}

\begin{definition}[Radix-Pass Arguments]
\lean{Light.RadixPass.Args}
\label{Light.RadixPass.Args}
\leanok
An argument record for a radix pass consists of natural numbers $w$, $perm$, $key$,
$stride$, $off$, $nb$, $cnt$, and $R$. Here $w$ is the number of indices; the index list
begins at cell $perm$, with the next $w$ cells reserved as scratch space; the key of
index $i$ is stored at cell $key+i\cdot stride+off$; every key is less than $nb$; the
$nb+1$ counters begin at cell $cnt$; and only cells at or beyond $perm$ among the first
$R$ cells may change.
\end{definition}

\begin{definition}[Counting-sort arguments for a radix pass]
\lean{Light.RadixPass.Args.sort}
\label{Light.RadixPass.Args.sort}
\leanok
\uses{Light.CountSort.Args, Light.RadixPass.Args}
A radix-pass argument record determines counting-sort arguments with item count $w$,
source address $\mathrm{perm}$, destination address $\mathrm{perm}+w$, and the same key
address, stride, offset, key bound, and counter address as the radix-pass record.
\end{definition}

\begin{definition}[Radix pass preconditions]
\lean{Light.RadixPass.Pre}
\label{Light.RadixPass.Pre}
\leanok
\uses{Light.Limits, Light.RadixPass.Args, Light.SegN}
For limits $lim$, memory $\mu$, radix-pass arguments $A$, a list of indices $\pi$, and a
key function $kd$, the precondition requires $lim.space\le lim.word$, $|\pi|=A.w$, and
every entry of $\pi$ is below $A.w$. The list is stored at $A.perm$. For each $i<A.w$,
the key cell $A.key+i\,A.stride+A.off$ lies before $A.perm$ and contains $kd(i)$, with
$kd(i)<A.nb$. The counter region starts after the index list and its $A.w$ scratch cells,
ends within the allowed region of length $A.R$ starting at $A.perm$, and that allowed
region lies within the space bound.
\end{definition}

\begin{definition}[Radix pass postconditions]
\lean{Light.RadixPass.Post}
\label{Light.RadixPass.Post}
\leanok
\uses{Light.IsStablePass, Light.RadixPass.Args, Light.SameOutside, Light.SegN, Light.cntLt, Light.keyAt}
Let $\mu:\mathbb{N}\to\mathbb{Z}$ be memory before the pass,
$\mu':\mathbb{N}\to\mathbb{Z}$ memory after it, $A$ be radix-pass arguments, $\pi$ a
list of indices, and $kd:\mathbb{N}\to\mathbb{N}$ their keys. The memories agree outside
the region of $A.R$ cells starting at $A.\mathrm{perm}$; for every $t\le A.\mathrm{nb}$,
the counter cell $A.\mathrm{cnt}+t$ contains the number of entries in the key sequence
of $\pi$ that are less than $t$. In addition, there is a list $\pi'$ that is a stable
ordering of $\pi$ by $kd$ and whose entries are stored consecutively in memory starting
at $A.\mathrm{perm}$.
\end{definition}

\begin{theorem}[Radix pass counting-sort preconditions]
\lean{Light.RadixPass.Pre.sort}
\label{Light.RadixPass.Pre.sort}
\leanok
\uses{Light.CountSort.Pre, Light.RadixPass.Args.sort, Light.RadixPass.Pre, Light.SegN.read, Light.keyAt}
\difficulty{4}
Let $L$ have word bound $L.word$, space bound $L.space$, and depth bound $L.depth$; let
$\mu:\mathbb{N}\to\mathbb{Z}$ be a memory map, let $A$ be radix-pass arguments, let
$\pi$ be a list of natural numbers, and let $kd:\mathbb{N}\to\mathbb{N}$. Suppose
$L.space\le L.word$, $|\pi|=A.w$, every entry of $\pi$ is less than $A.w$, and the
entries of $\pi$ are stored consecutively starting at $A.perm$. For every $i<A.w$,
suppose the key cell $A.key+i\cdot A.stride+A.off$ lies before $A.perm$, stores $kd(i)$,
and $kd(i)<A.nb$; also suppose $A.perm+2A.w\le A.cnt$, $A.cnt+(A.nb+1)\le A.perm+A.R$,
and $A.perm+A.R\le L.space$. Then the counting-sort arguments with $A.w$ items, source
address $A.perm$, destination address $A.perm+A.w$, and the same key parameters, key
bound, and counter address satisfy the counting-sort preconditions: item $j$ is the
entry of $\pi$ at position $j$, or $0$ if $j\ge|\pi|$, and its key is $kd$ of that item.
In particular, the source, destination, and counter regions fit in memory and are
pairwise disjoint; each item key cell fits in memory, lies outside the destination and
counter regions, and stores a key less than $A.nb$; and the destination and counter
regions fit in memory and are disjoint.
\end{theorem}

\begin{theorem}[Radix pass after copying sorted indices]
\lean{Light.RadixPass.Post.of_sorted_copied}
\label{Light.RadixPass.Post.of_sorted_copied}
\leanok
\uses{Light.CountSort.Post, Light.RadixPass.Args.sort, Light.RadixPass.Post, Light.RadixPass.Pre, Light.SameOutside, Light.SameOutside2, Light.exists_stablePass, Light.keyAt}
\difficulty{5}
Let $A$ be radix-pass arguments, let $\mu$ be the initial memory, let $\pi$ be a list of
natural-number indices, and let $kd$ assign a natural-number key to each index. Suppose
the radix-pass precondition holds, and the associated counting-sort postcondition holds
in memory $\mu_1$, with the item at position $j$ taken to be the $j$th entry of $\pi$
(or zero when $j$ is out of range) and its key taken to be the key of that item (using
zero for an out-of-range entry). Suppose also that $\mu_2(A.perm+i)=\mu_1(A.perm+A.w+i)$
for every $i<A.w$, and that $\mu_1$ and $\mu_2$ agree outside the interval starting at
$A.perm$ of length $A.w$. Then the radix-pass postcondition holds for $\mu_2$.
\end{theorem}

\begin{theorem}[Correctness and Stability of a Radix Pass]
\lean{Light.radixPass_ends}
\label{Light.radixPass_ends}
\leanok
\uses{Light.Ends, Light.Ends.callLast, Light.Ends.callThen, Light.Expr, Light.Expr.Gives, Light.Expr.cost, Light.RadixPass.Num, Light.RadixPass.Perm, Light.RadixPass.Post, Light.RadixPass.Post.of_sorted_copied, Light.RadixPass.Pre, Light.RadixPass.Pre.sort, Light.RadixPass.Unused, Light.copyBody, Light.copy_meets, Light.countSortBody, Light.countSort_meets, Light.radixPass}
\difficulty{5}
Let a resource-limited program have procedures at positions $pc$ and $pp$ containing the
counting-sort and copy routines, respectively. Let $A$ describe a radix pass on $w$
indices stored from address $perm$, with keys $kd(i)$ at addresses $key+i\,stride+off$,
all below $nb$; let $\pi$ be a permutation list of these indices, and let $\mu$ be the
initial memory. Assume $\pi$ has length $w$, all its entries are below $w$, it is stored
at $perm$, each key cell contains $kd(i)<nb$, the key cells precede $perm$, the counters
start at $cnt$ after the index and scratch regions, the counter region of length $nb+1$
lies within the $R$ cells from $perm$, and $perm+R$ is within the space limit; also
assume the space limit is at most the word bound. Suppose the input variables give $w$,
$perm$, $key$, $stride$, $off$, $nb$, and $cnt$ safely, the recursion depth is below the
depth limit, and $84w+58nb+72$ plus the costs of those five expressions is at most $T$.
Then the radix-pass program has an execution of cost at most $T$ after which only the
designated unused variable may have changed, all memory outside the $R$ cells from
$perm$ is unchanged, each counter at $cnt+t$ for $t\le nb$ equals the number of indices
$i<w$ with $kd(\pi_i)<t$, and the indices stored from $perm$ form a stable permutation
of $\pi$ ordered by nondecreasing $kd$ values.
\end{theorem}
